\section{Network Decisions}
The receiver gives the network evidence about each TB before decoding finishes. We use it for two decisions, summarized in Table~\ref{tab:network}: when to fall back, from the memory or to another receiver, and when to request a retransmission. As a memory entry only depends on the pilots of its slot, decoding a TB again does not change later slots.

\begin{table}[t]
\centering
\caption{Network decisions with 16-QAM over the three conditions, without (No PC) and with power control (PC); early retransmission and HARQ with power control. Stateless NRX: Large NRX retrained with $N_{\mathrm{it}}=2$.}
\label{tab:network}
\scriptsize\setlength{\tabcolsep}{3.5pt}
% written by scripts/make_nrx_paper_tables.py; do not edit
\begin{tabular}{@{}lrr@{}}
\toprule
 & No PC & PC \\
\midrule
\multicolumn{3}{@{}l}{\emph{Fallback from the memory}: DMRS symbols on which} \\
\multicolumn{3}{@{}l}{the channel filter uses no stored TB [\%]} \\
\quad at a channel change & -- & 98.4 \\
\quad after a phase jump & -- & 93.4 \\
\quad otherwise, at most & -- & 0.6 \\
\midrule
\multicolumn{3}{@{}l}{\emph{Fallback after a CRC failure}: failed TBs rescued [\%] by} \\
\quad LS + LMMSE & 0.0 & 0.0 \\
\quad Stateless NRX, $N_{\mathrm{it}}=2$ & 0.0 & 0.2 \\
\quad LMMSE + K-Best & 5.9 & 4.7 \\
\quad Large NRX, $N_{\mathrm{it}}=8$ & 3.3 & 7.8 \\
Failed TBs after an error of the UE [\%] & 97.1 & 80.9 \\
\midrule
\emph{Early retransmission} & CRC hist. & Controller \\
Errors caught early [\%] & & \\
\quad proposed NRX & 79.5 & 97.3 \\
\quad proposed NRX, no memory state & 79.5 & 96.9 \\
\quad stateless NRX, no memory state & 82.1 & 90.3 \\
Correct TBs flagged [\%], proposed & 2.7 & 2.7 \\
Errors after the CRC per 100 TBs & & \\
\quad proposed NRX, no memory state & 2.4 & 0.4 \\
\quad stateless NRX, no memory state & 4.0 & 2.2 \\
\midrule
\emph{Simulated HARQ} & After CRC & Early \\
Median time to decoding [ms] & 3.9 & 3.4 \\
Unneeded transmissions per 100 TBs & 0.0 & 2.4 \\
\bottomrule
\end{tabular}

\end{table}

\subsection{Fallback}
The receiver itself decides when to fall back from its memory: on the current pilots, the channel filter chooses how many stored TBs to use, down to none. With power control, it uses none at almost every channel change or phase jump, and almost never otherwise, a decision the classical multi-slot estimator cannot make.

After a CRC failure, a fallback that a real-time BS could run, LS + LMMSE or a stateless NRX with $N_{\mathrm{it}}=2$, rescues almost none of the failed TBs, and LMMSE + K-Best and the Large NRX with $N_{\mathrm{it}}=8$ rescue only a few. Most failed TBs follow an error of the same UE, so the remedy is to raise that UE's power, as power control does, not to switch the receiver.

\subsection{Early Retransmission}
Requesting a hybrid automatic repeat request (HARQ) retransmission before decoding finishes, known as early HARQ, shortens the recovery from an error~\cite{goektepe2018subcode,strodthoff2019early}. Our controller checks each TB three times before its decoding finishes: once the CRC of the UE's previous TB is known, after detection, and after the first eight belief propagation iterations. At each check, a logistic regression model maps the evidence available so far to a failure score: first, the UE's CRC history, the noise level, and the power of the LS and refined estimates and the memory state; then, the bit error rate implied by the NRX LLRs, i.e., the mean of $1/(1+e^{|\ell_{u,j}|})$ over the TB; and finally, the unsatisfied parity checks and the bit error rate of the posterior LLRs. A retransmission is requested at the first check whose score exceeds its threshold; a missed error is still retransmitted after the CRC fails, so the controller cannot impair decoding. We fit the models and thresholds on 16 sequences of each condition and test them on the other 16, with power control and at the SNR point closest to a TBLER of 10\%.

Requesting a retransmission whenever the UE's previous TB failed already catches most errors in the test sequences. With thresholds that flag the same share of correct TBs in the fitting sequences, the controller catches almost all of them, and flags as many correct TBs in the test sequences as this rule. Without the memory state, it still catches almost all errors of the proposed NRX, and clearly more errors of the stateless NRX than the rule. As that NRX fails more often at the same SNR (Fig.~\ref{fig:tbler}(d)--(f)), many more errors remain for retransmission after the CRC fails than with the proposed NRX. Simulating the retransmissions with Chase combining shows that early requests let the BS decode a failed TB one slot earlier in median, at the cost of a few unneeded transmissions.
